% =========================================================
% PAGE 6: REPORT DETAILS / METHODOLOGY (ARABIC)
% =========================================================

% 1. Reset Geometry (Standard Internal Layout)
\newgeometry{left=2cm, right=2cm, top=1cm, bottom=2.5cm}



% ---------------------------------------------------------
% TITLE SECTION
% ---------------------------------------------------------
\noindent
{\fontsize{24}{28}\selectfont \textbf{\color{ecestitle} التقرير التفصيلي}}
\vspace{0.3em}
\label{sec:ar_detailed_report}
\noindent
\\\\
{\fontsize{18}{22}\selectfont \textbf{\color{ecesteal} المنهجية}}
\label{sec:ar_methodology}
% ---------------------------------------------------------
% PAGE 6 CONTENT
% ---------------------------------------------------------
\noindent
\\\\
{\fontsize{11}{15}\selectfont \color{textdark}
استكمالا لجهود المركز المصري للدراسات الاقتصادية في توفير معلومات متكاملة تعكس التطورات التي يشهدها الاقتصاد المصري بشكل عام ومجتمع الأعمال بشكل خاص يقوم المركز بإصدار \textbf{بارومتر الأعمال} منذ 1998. يقدم البارومتر تقييما دوريا ربع سنوي لأداء عينة من شركات القطاع الخاص تغطي مختلف القطاعات والأحجام. ويعكس هذا التقييم رأي مجتمع الأعمال بشأن التطورات التي شهدتها مجموعة من المتغيرات وذلك خلال الربع محل الدراسة كما يلقي الضوء على توقعاتها لتطورات نفس مجموعة المتغيرات خلال الربع التالي.
}

\vspace{0.5em}

% ---------------------------------------------------------
% DIAGRAM 1: Indicators (ch4.png)
% Made smaller (Height max 3.5cm) to save space
% ---------------------------------------------------------
\begin{figure}[h!]
    \centering
    \includegraphics[width=0.95\linewidth, height=3.5cm, keepaspectratio]{ch4.png}
\end{figure}

\vspace{0.2em}

\noindent
{\fontsize{11}{15}\selectfont \color{textdark}
ويزيد من أهمية هذا العدد من بارومتر الأعمال ما يواجه مجتمع الأعمال منذ مطلع عام 2020 من تحديات نتيجة لجائحة كورونا تبعها تعاف مصحوب بالعديد من التحديات خلال عام 2021 ثم الحرب الروسية الأوكرانية في مطلع عام 2022. بالإضافة إلى الاضطرابات الجيوسياسية في الشرق الأوسط والبحر الأحمر منذ أكتوبر 2023، مما فاقم من هذه التحديات؛ لذا من المهم تتبع تأثير هذه المستجدات عليه.

\vspace{0.5em}

ويستعرض هذا التقرير تقييما لأداء شركات العينة خلال الربع \textbf{(\CurrentQuarterText)} وتوقعاتها للربع \textbf{(\NextQuarterText)}.

\vspace{0.5em}

يبدأ التقرير بإلقاء نظرة عامة على الاقتصاد الكلي على المستويين العالمي والمحلي، ثم يقوم باستعراض نتائج تقييم الأداء والتوقعات على مستوى المؤشر الإجمالي، بعدها ينتقل إلى المعوقات التي واجهت مجتمع الأعمال خلال الربع محل الدراسة، وأولويات تحسين مناخ الأعمال من وجهة نظر شركات العينة. وأخيرا، يختتم التقرير بتقييم الأداء والتوقعات على مستوى المؤشرات الفرعية.

\vspace{0.5em}

وقد تم بناء بارومتر الأعمال استنادا إلى نتائج مسح يجريه المركز بصفة دورية كل ثلاثة شهور لعينة ثابتة تضم \textbf{120 شركة} من شركات القطاع الخاص موزعة على النحو التالي:
}

\vspace{0.2em}

% ---------------------------------------------------------
% DIAGRAM 2: Sample Distribution (ch5.png)
% Made smaller (Height max 4.5cm) to ensure it fits on Page 6
% ---------------------------------------------------------
\begin{figure}[h!]
    \centering
    \includegraphics[width=0.85\linewidth, height=4.5cm, keepaspectratio]{ch5.png}
\end{figure}

\clearpage

% =========================================================
% PAGE 7: METHODOLOGY CONTINUED (ARABIC)
% =========================================================

% Geometry Persists from Page 6


% ---------------------------------------------------------
% METHODOLOGY BULLET POINTS
% ---------------------------------------------------------
\vspace*{0.5cm} % Small top padding

\begin{itemize} \setlength\itemsep{0.6em}
    \item \fontsize{11}{15}\selectfont \color{textdark} يقوم التحليل بتقييم أداء شركات العينة خلال الربع محل الدراسة وتوقعاتها للربع التالي مع مقارنتها بنتائج الربع السابق والربع المناظر من العام السابق.
    
    \item \fontsize{11}{15}\selectfont يتم تقييم الأداء والتوقعات على مستويين هما: نتائج المؤشر الإجمالي، ونتائج المؤشرات الفرعية.
    
    \item \fontsize{11}{15}\selectfont يمثل مؤشر بارومتر الأعمال متوسطا بسيطا لمجموعة المؤشرات الفرعية للمتغيرات الواردة في الاستبيان. ويأخذ المؤشر الإجمالي قيما أكبر أو أقل أو مساوية للمستوى المحايد (50 نقطة).
\end{itemize}

\vspace{0.5em}

% ---------------------------------------------------------
% CENTRAL GRAPHIC (Gauge / ch6.png)
% Reduced size for balance
% ---------------------------------------------------------
\begin{figure}[h!]
    \centering
    \includegraphics[width=0.6\linewidth, height=4cm, keepaspectratio]{ch6.png}
\end{figure}

\vspace{1em}

% ---------------------------------------------------------
% EQUATION SECTION (Styled Box)
% ---------------------------------------------------------
\noindent
\begin{tikzpicture}
    \setRTL
    \node[inner sep=0pt] (container) {
        \begin{minipage}{\textwidth}
            \centering
            % Text
            {\fontsize{11}{15}\selectfont \color{textdark} ويتم حساب المؤشرات الفرعية لكل متغير باستخدام المعادلة الآتية:}
            \vspace{0.8em}
            
            % Equation Box
            \colorbox{ecesteal}{
                \begin{minipage}{5cm}
                    \centering
                    \vspace{0.3cm}
                    \color{white} \large \bfseries
                    $X = \displaystyle \frac{I + S}{100 + S} \times 100$
                    \vspace{0.3cm}
                \end{minipage}
            }
            \vspace{0.8em}
            
            % Explanation
            {\fontsize{10}{13}\selectfont \color{textdark} 
            حيث $I$ هي نسبة الشركات التي أفادت بزيادة المتغير، و$S$ هي نسبة الشركات التي أفادت بثباته.
            }
        \end{minipage}
    };
\end{tikzpicture}

\vspace{1.5em}

% ---------------------------------------------------------
% CHALLENGES & PRIORITIES TEXT
% ---------------------------------------------------------
\noindent
{\fontsize{12}{15}\selectfont \textbf{\color{ecestitle} فيما يتعلق بالمعوقات وأولويات تحسين مناخ الأعمال:}}

\vspace{0.8em}

\noindent
{\fontsize{11}{15}\selectfont \color{textdark}
تقوم الشركات بتقييم شدة كل معوق بحيث يتراوح التقييم من صفر (معوق غير مؤثر)، إلى 4 (معوق شديد التأثير). ويُسمح للشركة باختيار أكثر من معوق واحد. وفيما يتعلق بأولويات تحسين مناخ الأعمال، تقوم الشركات بتقييم أولويات تحسين مناخ الأعمال بحيث يتراوح تقييم كل محور من صفر (محور غير هام)، إلى 4 (محور له أولوية كبيرة). ويُسمح للشركات باختيار أكثر من محور كأولوية لتحسين مناخ الأعمال.

\vspace{0.8em}

يلي ذلك حساب متوسط مرجح لأعداد الشركات وتقييمها للمعوق/ الأولوية على مستوى العينة ككل.

\vspace{0.8em}

يُعاد تقييم متوسطات كافة المعوقات/ الأولويات لتتراوح ما بين صفر و1 ثم يتم عمل معايرة Normalization على مستوى القيم الجديدة لمتوسطات كافة المعوقات/ الأولويات بحيث يمكن ترتيب المعوقات/ الأولويات تنازليا حسب شدتهم فتكون النسبة 100 أشد معوق وأكبر أولوية.
}

\clearpage